% ============================================
% DOC-20 (EN)
% English translation — SPEC-935-R355
% ============================================
\section*{\textbf{DOC-20} --- Chain of Custody Record}
\markboth{DOC-20 --- Chain of Custody Record}{}

\noindent \textit{Document extracted from the forensic file.}

\vspace{0.5em}\noindent\rule{\textwidth}{0.5pt}\vspace{0.5em}

\noindent {\footnotesize\texttt{PUBLIC ARCHIVE $\cdot$ DOCUMENTARY EVIDENCE $\cdot$ CASE 1,260}}\\[-0.2em]

\noindent \textbf{Forensic Medical Institute of Barbacena --- Control Document}\\
\noindent \textbf{Object:} leather-bound notebook, dimensions 16\,cm × 22\,cm, approximately 200 pages, without original numbering. \\
\noindent \textbf{Attributed to:} Joaquim Antônio Correia (Patient 1,260). \\
\noindent \textbf{Nature:} personal diary / unofficial clinical record.

\underline{Record of the object's movement:}

\begin{moltablefive}
\hline
\noindent \textbf{No.} & \textbf{Date} & \textbf{Origin} & \textbf{Destination} & \textbf{Responsible Party} \\ \hline
\noindent 1 & 14/Jun/1979 & Grave 1,260, Colony Cemetery & IML --- Autopsy & Genésio Barbosa (gravedigger) \\
\noindent 2 & 14/Jun/1979 & IML --- Autopsy & IML --- Vault & Dr. Antônio Tavares (forensic pathologist) \\
\noindent 3 & 16/Jun/1979 & IML --- Vault & \textit{[unfilled]} & Dr. Heitor Oliveira \\
\noindent 4 & \textit{[illegible]} & \textit{[not recorded]} & Minas Gerais Public Archive & \textit{[unidentified]} \\
\noindent 5 & 12/Jan/2026 & MGP Archive --- Consultation & MGP Archive --- Research & Dr. Lúcia Mendes \\
\noindent 6 & 28/Jul/2026 & MGP Archive --- Research & MGP Archive --- Vault & Dr. Lúcia Mendes \\
\noindent 7 & \textit{[blank]} & \textit{[blank]} & \textit{[blank]} & \textit{[blank]} \\
\hline
\end{moltablefive}

\noindent \textbf{Attached observations:}

\noindent \textit{1. Entry No. 3 (Dr. Heitor Oliveira) has no formal removal authorization. The IML Register Book records the physician's name, but there is no witness signature, stamp of competent authority, or written justification.}
\noindent \textit{2. Entry No. 4 was reconstructed from indirect mention in internal correspondence of the Minas Gerais Public Archive. The original date and the name of the person responsible for the transfer were removed from the physical document, leaving only an empty rectangle where text used to be. Optical analysis suggests the name was scraped off with a blade.}
\noindent 3. Field No. 7 was left blank. In 2026, when the object was transferred to the definitive vault, the employee responsible for filling it out --- whose name does not appear in the document --- stated verbally that ``it was not up to him to fill in the last line''. Asked why, he answered: \textquotedblleft{}\textit{Because the diary is not going to stay here.}\textquotedblright{}

\noindent 4. When asked if he knew where the object would be taken next, the employee hesitated and said: \textquotedblleft{}\textit{Into the hands of whoever is waiting.}\textquotedblright{}

\noindent \textbf{Archivist's conclusion:}
\noindent The object's chain of custody is not continuous. There are two unfilled documentary gaps and one entry whose responsible party cannot be confirmed. The object circulated for at least 47 years without a formal record of its whereabouts.
\noindent This record is, itself, a document of absences.

\noindent \textit{Archivist's note: the fragments of this part were organized by Dr. Lúcia Mendes during her research in 2026. Her annotations are in italics in the margins of the original documents. The oldest document in the archive --- from 1853 --- was found by her and attached to the case file.}

\vspace{0.5em}
\noindent\textit{\textbf{↪ Routes:} \rota{DOC-01} \rotasep \rota{LUC-02}}
